\newcommand{\figuraInsectoTensionSuperficial}{%
\begin{figure}[H]
\centering
\begin{tikzpicture}[>=Latex,font=\small,line cap=round,line join=round]
  \fill[blue!10] (0,0) rectangle (12,2.45);
  \draw[very thick,blue!65!black] (0,2.45)
    .. controls (1.2,2.45) and (1.45,2.10) .. (2.1,2.35)
    .. controls (3.0,2.65) and (3.25,2.10) .. (4.0,2.35)
    .. controls (5.0,2.65) and (5.35,2.05) .. (6.0,2.35)
    .. controls (7.0,2.65) and (7.35,2.05) .. (8.0,2.35)
    .. controls (9.0,2.65) and (9.3,2.10) .. (10.0,2.35)
    .. controls (10.8,2.60) and (11.3,2.45) .. (12,2.45);
  \fill[black!75] (6,4.15) ellipse (1.05 and 0.38);
  \fill[black!70] (7.1,4.15) circle (0.34);
  \draw[very thick] (5.45,3.95)--(4.0,2.35);
  \draw[very thick] (5.85,3.90)--(6.0,2.35);
  \draw[very thick] (6.25,3.90)--(8.0,2.35);
  \draw[very thick] (6.65,4.00)--(10.0,2.35);
  \draw[very thick] (5.35,4.15)--(2.1,2.35);
  \draw[very thick] (6.85,4.25)--(9.65,3.05);
  \draw[->,ultra thick,purple!75!black] (6,4.55)--(6,3.05)
    node[midway,right] {$W$};
  \foreach \x in {2.1,4,6,8,10}
    \draw[->,thick,red!75!black] (\x,2.18)--(\x,2.82);
  \node[red!75!black,align=center] at (3.0,1.15)
    {componentes verticales\\[1mm]de la tensión superficial};
  \node[blue!65!black] at (9.8,0.75) {agua};
\end{tikzpicture}
\caption{Un insecto zapatero deforma la superficie sin romperla. Las
componentes verticales de la tensión superficial alrededor de cada pata
contribuyen a sostener su peso.}
\label{fig:insecto-tension-superficial}
\end{figure}%
}